%
% Copyright 2014, NICTA
%
% This software may be distributed and modified according to the terms of
% the BSD 2-Clause license. Note that NO WARRANTY is provided.
% See "LICENSE_BSD2.txt" for details.
%
% @TAG(NICTA_BSD)
%

%%%%%%%%%%%%%%%%%%%%%%%%%%%%%%%%%%%%%%%%%%%%%%%%%%%%%%%%%%%% 
% For papers:
% - use ``article''
% - remove the \else part of the \iftrue below (and remove the
%    \iftrue/\fi commands)
%    \documentclass[a4paper,11pt,twoside,twocolumn]{article}
%%%%%%%%%%%%%%%%%%%%%%%%%%%%%%%%%%%%%%%%%%%%%%%%%%%%%%%%%%%%    
%    For technical reports, manuals etc:
%    - use ``report''
%    - use the \else part of the \iftrue below (and remove the \iftrue/fi commands
%    Use ``report'' for techreports, manuals etc and use \else part of \iffalse
%    and also uncomment the \usepackage{disy}

\documentclass[a4paper,11pt,twoside]{report}
\usepackage[colour,nictaonly]{disy}

%%%%%%%%%%%%%%%%%%%%%%%%%%%%%%%%%%%%%%%%%%%%%%%%%%%%%%%%%%%%
\usepackage{msctexen}

%%% $Id: paper.tex,v 1.22 2008-07-19 08:14:55 gernot Exp $

% Setting this to true turns on the ``draft'' watermark
\newif \ifDraft         \Drafttrue

%%% Hyperlinks embed links for references in the PDF output (both for
%%% internal as well as external links).
\newif \ifhyperlinks    \hyperlinkstrue
%%% When using this you should use \autoref{label} instead of
%%% Section~\ref{label}. This will make ``Section'' part of the link, not
%%% just the number.

% Easy control over page margins. You can use ``right'', ``top'', ...
\usepackage[margin=25mm]{geometry}
% \usepackage[margin=25mm,left=33mm]{geometry}	% better for double-sided reports
\usepackage{xspace} 
\usepackage{graphicx}
\setkeys{Gin}{keepaspectratio=true,clip=true,draft=false,width=\linewidth}
\usepackage{times,cite,url,fancyhdr}
\urlstyle{rm}
\renewcommand{\ttdefault}{cmtt}	% CM rather than courier for \tt

\ifDraft
\usepackage{draftcopy}
\newcommand{\Comment}[1]{\textbf{\textsl{#1}}}
\newcommand{\FIXME}[1]{\textbf{\textsl{FIXME: #1}}}
\date{\small$\relax$Id: paper.tex,v 1.22 2008-07-19 08:14:55 gernot Exp $\relax$}
\else
\newcommand{\Comment}[1]{\relax}
\newcommand{\FIXME}[1]{\relax}
\date{}
\fi

\ifhyperlinks
\usepackage[pdftex,pagebackref,hyperindex,bookmarks]{hyperref}
\hypersetup{linktocpage,colorlinks}
\else   
\providecommand{\href}[2]{\url{#2}}
\fi

\pagestyle{fancyplain}
\lhead[\fancyplain{}{\sl\thepage}]{\fancyplain{}{\sl\rightmark}}
\chead{}
\rhead[\fancyplain{}{\sl\leftmark}]{\fancyplain{}{\sl\thepage}}
\lfoot[\fancyplain{\sl\thepage}{}]{}
\cfoot{}
\rfoot[]{\fancyplain{\sl\thepage}{}}

\newcommand{\ty}[1]{\texttt{#1}}

\begin{document}

\ifhyperlinks

%% Capitalisation of cross references
\renewcommand{\chapterautorefname}{Chapter}
\renewcommand{\sectionautorefname}{Section}
\renewcommand{\subsectionautorefname}{Section}
\renewcommand{\subsubsectionautorefname}{Section}
\renewcommand{\appendixautorefname}{Appendix}
\renewcommand{\Hfootnoteautorefname}{Footnote}

%% Commands for index
\newcommand{\Htextbf}[1]{\textbf{\hyperpage{#1}}}
\fi

%% Custom Colours
\definecolor{refos_purple}{RGB}{163, 73, 164}
\definecolor{refos_dpurple}{RGB}{72, 12, 68}
\definecolor{refos_dgray}{RGB}{64, 64, 64}
\definecolor{seL4_green}{RGB}{0, 64, 0}

%% Custom commands
\newcommand{\refOS}{\textcolor{refos_purple}{\texttt{RefOS}}\xspace}
\newcommand{\Lf}{\texttt{L4}\xspace}
\newcommand{\seLf}{\textcolor{seL4_green}{\texttt{seL4}}\xspace}

\title {
    \centering\includegraphics[width=0.5\textwidth]{refos} \\\
    RefOS --- Reference Design For A Microkernel Based Operating System
}

\author {Kevin Elphinstone, Andi Cheng and Xi Chen}
\AuthorEmail {
    Kevin.Elphinstone@nicta.com.au,
    Andi.Cheng@nicta.com.au,
    XiMa.Chen@nicta.com.au
}
\date{\today}

\begin{titlepage}
\maketitle
\end{titlepage}

\urlstyle{sf}
\thispagestyle{empty}

\begin{abstract}
  \pagenumbering{roman}

  The design of a multi-server operating system is crucial to its efficiency, security and
  dependability. Increasing usage of embedded systems and microcontroller chips has increased the
  interest from both industrial and academic communities in microkernel-based component-based
  software systems. This has resulted in heightened interest in developing better systems and more
  advanced microkernels.
  
  This document outlines a basic reference design for system components and interfaces for an
  operating system. The protocols have been designed with an \Lf like microkernel in mind, assuming
  functionality such as inter-thread IPC messaging, capability-based access models, virtual memory
  support, and so forth. The reference design aims towards standardising the protocols for tasks
  often required to build a system such as sharing, virtual memory management, and naming. Operating
  systems may implement these protocols and / or extend it to incorporate more features.

  A sample implementation of the protocol is provided, in the form of a simple OS able to run simple
  client programs. The sample implementation uses the \seLf microkernel, adding an extra dimension
  of dependability.

  \vfill
  
  \noindent
  \copyright~2014 NICTA
\end{abstract}
\setcounter{page}{1}


\tableofcontents

\setcounter{page}{1}
\pagenumbering{arabic}

\newcommand{\srv}[1]{\textcolor{refos_dpurple}{\texttt{#1}}}
\newcommand{\obj}[1]{\texttt{#1}}
\newcommand{\cp}[2]{\textcolor{refos_dpurple}{\srv{#1}\texttt{\_}\obj{#2}\texttt{\_C}}}
\newcommand{\met}[3]{\cp{#1}{#2}{\texttt{.#3}}}
\newcommand{\pro}[4]{
    \met{#1}{#2}{#3} \\
    $\Rightarrow$ \textcolor{refos_dgray}{\texttt{#4}} \\
}

\chapter{Introduction}
\label{s:intro}
%
% Copyright 2014, NICTA
%
% This software may be distributed and modified according to the terms of
% the BSD 2-Clause license. Note that NO WARRANTY is provided.
% See "LICENSE_BSD2.txt" for details.
%
% @TAG(NICTA_BSD)
%

This document provides the abstract protocols for a multi-server microkernel based operating system.

The overall design revolves around the abstraction of a dataspace. A dataspace is a memory space, a
series of bytes, representing anything from physical RAM to file contents on disk to a device or
even to a random number generator. The concept is analogous to UNIX files, which may represent
anything from \texttt{/usr/bin/sh} to \texttt{/dev/audio} to \texttt{/dev/urandom}.

This dataspace abstraction which dataspace servers provide paves the way for more complicated
infrastructure on top such as sharing with complex trust relationships, memory management, file
systems, and distributed naming. On top of this additional infrastructure, the OS may implement more
features such as POSIX standard syscalls, ports of existing software, drivers, display servers, and
so forth.

% ----------------------------------------------------------------------
\section{Terms}

This section will attempt to briefly clarify term ambiguities before the details
are presented. The brief descriptions also quickly give an overview of the
terms.

\texttt{Methods} are conceptual "function calls", although they are usually
implemented via communication with another process, in which case they are
actually remote procedure calls. Regardless of the underlying implementation,
methods here refer to the actual procedures of an interface which get invoked.

\texttt{Interface} is a collection of (usually related) methods, that
usually serve to abstract a conceptual object. For instance, the \texttt{c}
standard library functions \texttt{fopen}, \texttt{fscanf}, \texttt{fprintf},
\texttt{fclose} ..etc form an abstract interface and provide abstraction for a
\emph{file} object.

\texttt{Message Sequence} is an ordered sequence of method invocations,
replies, event notifications or operations which may achieve some overall
result or goal.  Example message sequences using the described protocols for
commonly needed goals will be presented in this document.

\texttt{Protocol} here will refer to a collection of interfaces, their methods,
message sequences and system components.

\texttt{Dataspaces} are abstractions over memory objects, analogous to UNIX files. A dataspace
may be opened / closed, written to or read from, mapped into client memory, and may represent
a device, a file, or simply anonymous memory.

\texttt{Servers} are normal processes which have registered themselves using
the naming service so that other client processes may find. Once the client
finds the process, the client may then connect and establish a session with the
server, and the server is given access to the notification of client deaths in
order to do any client book-keeping it may require.

\texttt{Dataspace servers} (henceforth also referred to as \emph{dataservers}) are
servers which implement the dataspace interface (or a subset of it), in order to
provide a dataspace service to clients. Dataspace servers may be most likely
designed to provide the dataspaces which represent a common concept. For
instance, a file server may implement the dataspace interface for files stored
on disk, while an audio driver may serve audio bytestream dataspaces.

\texttt{Process servers} are servers which implement the process server interface,
and provide process abstraction. In practice, the process server may also need
to implement additional interfaces and provide additional functionality in order
for a system to start. For instance in our sample implementation the process server
also does naming, console input/output, and implements the dataspace interface
for anonymous memory and acts as the memory manager.

\texttt{Pagers} are servers which implement the pager service interface, which
is very closely related to the dataspace interface, and can be considered an
optional feature of a dataspace server. Dataspace servers which act as pagers
may have its dataspaces be mapped directly into client virtual memory.

% ----------------------------------------------------------------------
\section{Microkernel}

The \refOS protocol is designed with an \Lf family microkernel in mind, and
assumes the existence of features such as inter-thread IPC and capabilities.
This section describes in detail the notable assumed features, with brief
descriptions of such features.  These features are available in advanced
\Lf-based microkernels such as \seLf.

\subsection{IPC Messages}

\begin{itemize}

	\item 
	\textbf{Synchronous Calls} are used to invoke methods implemented by a
	server. They involve a badged cap invocation (of a synchronous endpoint or
	SEP) that identifies the server and represents the authority of the caller.
	The caller blocks until the server responds. \emph{Note: A caller
	implicitly trusts the server it calls.} Calling an untrusted server is
	possible, but should be avoided in general -- which we strive to do in
	\refOS.

	Exceptions, such as page faults, are propagated via sending and blocking on
	a SEP.

	\item 
	\textbf{Reply} via reply caps. A call via invoking a SEP generates a Reply
	capability for the server to respond to the caller. Thus the callee can
	send a response to the original caller, and unblock the caller, by replying
	with the reply capability. Note that reply caps are guaranteed to either
	succeed or fail and not block. This is required for a trusted server to
	safely reply to an untrusted client.

	\item
	\textbf{Asynchonous Notifications} are sent via a \texttt{Send} to an
	asynchronous end point (AEP) capability. Notifications are non-blocking
	with at-least-once semantics (the bits are OR'd together). This AEP can be
	used to notify an untrusted client. To send messages(as opposed to just
	notifications) to an untrusted client, shared memory must be set up between
	the two. Care must be taken to avoid blocking, denial of service, etc... on
	and shared memory buffer which may be filled up.

\end{itemize}

\subsection{Capability Access Model}

\refOS assumes the microkernel provides support for a capability-based access
model, and more importantly, the ability to transfer badged capabilities via
IPC. 

The basic concept of a capability (henceforth referred to as \emph{cap}) is that
an object representing the permission to access a particular object. This is
similar in idea to UNIX file descriptors, which represent the access that the
process has to the file. Just like a UNIX file descriptor table, the capability
objects are to be kept in kernel memory. User programs only have handles to
such objects, and cannot directly manipulate them.

The \refOS is designed with the assumption that synchronous call and reply IPC
messages may contain capabilities, and that there is some way to store a
minimal amount of immutable information in such capabilities, and that the
receiving process may read the immutable information from the capabilities that
it has received via IPC. 

In the sample implementation given using the seL4 microkernel, this is
implemented as kernel IPC endpoint \emph{badges}; endpoint capabilities may be
badged with an integer, and the badge of an endpoint once it has been badged is
immutable. When a badged endpoint cap is sent along the same endpoint the one
it is pointing to, then the kernel automatically \emph{unwraps} the badge to be
read. Therefore, only the process which created the badge can read it back via
IPC. \refOS makes heavily use of these badges in order for servers to track
server objects (such as windows and dataspaces).

\subsection{Virtual Memory Support}

\refOS assumes that the underlying kernel exposes virtual memory management.
This is a standard feature of microkernels. It is also assumed that the kernel
is able to deliver page fault exceptions via IPC to a memory manager thread,
and be able to resolve the fault at a later stage by replying to the IPC.

% -----------------------------------------------------------------------------
\section{Notation}

This section outlines the notation used to describe the \refOS protocol.
Specifically, it looks at the messaging between components that is used to
develop higher-level protocols. 

\subsection{Operation Types} 

\autoref{fig:mta} presents the four basic operations we are using (and
indirectly attempting to limit ourselves to). They are synchronous calls (i.e.
a method invocation on an object), replies, asynchronous notifications, and
kernel system calls that asynchronously modify the state of another process
(such as mapping memory directly into the virtual address space of another
process).

The synchronous / asynchronous calls and reply operations may be implemented
directly using the microkernel features described above, such as IPC,
capabilities, the ability to send capabilities via IPC, virtual memory
operations, and so on.

\begin{figure}[htb]
  \centering
  \begin{msc}
    msc {
      A, B;
      A => B [ label = "A synchronous call or fault via a SEP"];
      A << B [ label = "A return from a call (via reply cap)"];
      A -> B [ label = "An asynchronous signal via an AEP"];
      A :> B [ label = "A kernel operation that affects B"];
    }
  \end{msc}
  \caption{Message Type Arrows.}
  \label{fig:mta}
\end{figure}

\subsection{Message Notation} 

Protocols are shown as series of messages. Messages are implemented by operations,
and the types of different messages directly correlate to the operation types.
This section outlines the notation which this document will henceforth adhere to
when describing the various forms of messages that form the \refOS protocol.

\subsubsection*{Method Invocation}

Most of the interface methods are synchronous. The client making the call is
considered to be less trusted than the server that is receiving the call. As
the client is making the call, it is blocked until the server finished handling
the call and replies via the reply operation. When the server replies, the call
is said to be finished and the client resumes execution. Method invocation is
implemented with the synchronous call operation via a SEP. Capabilities may be
passed via method invocation, and also via the corresponding reply.

For synchronous method invocations by a less trusted client to a more trusted
server, the following notation is used:

\begin{center}
  \pro{server}{interface}{method(arg1, arg2, ...)}{(val1, val2, val3 ...)}
\end{center}

Where:

\begin{itemize} 
	\item
	\srv{server} indicates the name of the server that is receiving and
	handling the method invocation via an endpoint. Eg. \srv{procserv} and
	\srv{dataserv} will henceforth be used in method descriptions to referr to
	a process server and a data server respectively.
	
	\item
	\obj{interface} refers to the name of the interface that the server
	implements.

	\item
	\texttt{method(arg1, arg2, ...)} refers to the name of the interface
	method, \texttt{(arg1, arg2, ...)} representing the arguments of the method
	call.

	\item
	\texttt{(val1, val2, val3 ...)} representing the set of return values, output variables
	and/or reply capabilities of the method invocation.
\end{itemize}

For example, the notation \met{dataserv}{dataspace}{open} means an invocation
of the \texttt{open} method of the \obj{dataspace} interface on the server
called \srv{dataserv}. Note that any \srv{server} can implement a particular
\obj{interface}.  The encoding of method and arguments in the actual message is
left as an implementation detail.

\autoref{fig:epm} describes an example protocol description (in the form of a
message sequence diagram) of a client process invoking a method via an
synchronous message operation:

\begin{figure}[htb]
  \centering
  \begin{msc}
    msc {
      R [ label ="Server"],
      C [ label = "Client"];
      C => R [ label = "server\_process\_C.method(13, id)"];
    }
  \end{msc}
  \caption{Example protocol message sequence diagram.}
  \label{fig:epm}
\end{figure}

\subsubsection*{Event Notification}   

Some servers need a method of notifying another process without blocking and
waiting for the other process to finish processing it. Asynchronous signal
operations via asynchronous endpoints are used to send such notifications to
other servers. Capabilities may not be transferred via a notification.

Note that \refOS does not make the assumption that the kernel is able to pass
message contents via the notification itself, just that the kernel is able
to pass just the notification. In order for the actual message to be delivered
along with the notification, the implementation of the protocol may for instance
make use of a shared buffer.

The following notation is used for event notifications:

\begin{center}
\met{server}{event}{notify(arg1, ...)}
\end{center}

This notation is purposely very similar to a method invocation, as those
two are similar in concept in that they are both inter-process messages.
The different message sequence diagram arrows used for synch calls
and async notifications clearly distinguish between the two.

\subsubsection*{Abstracted protocols}

Sometimes the protocols are of different level of abstractions. In order to
separate two closely related protocols in one's sequence diagram, we use the
following representation as indicated in \autoref{fig:separate_p}:

\begin{figure}[htb]
  \centering
  \begin{msc}
    msc {
      A [ label ="Process A"],
      B [ label = "Process B"];
      A abox B [ label = "A separate protocol goes here."];
      A => B [ label = "processB\_object\_C.method(arg, ...)"];
    }
  \end{msc}
  \caption{Example of separation of protocols.}
  \label{fig:separate_p}
\end{figure}

%%% Local Variables: 
%%% mode: latex
%%% TeX-master: t
%%% End: 


\chapter{Interfaces}
\label{s:interface}
%
% Copyright 2014, NICTA
%
% This software may be distributed and modified according to the terms of
% the BSD 2-Clause license. Note that NO WARRANTY is provided.
% See "LICENSE_BSD2.txt" for details.
%
% @TAG(NICTA_BSD)
%

This chapter will provide the interfaces which make up the \refOS protocol. These interfaces will
set the language with which components in the system will then interact. An interface will also aim
to provide the abstraction to manage conceptual objects, such as a process or a file or a device.

Implementations may choose to extend these interfaces to provide any extra system functionality, as
well as attach additional interfaces.

% ----------------------------------------------------------------------
\section{Objects}

Servers implement interfaces, which provide management (i.e. creation,
destruction, monitoring, sharing, manipulation) object(s). For example, an
audio device may implement the dataspace interface for UNIX \texttt{/dev/audio},
which gives the abstraction to manage "audio" objects.

In our sample implementation, the access to an object (i.e. the permission to
invoke the methods which manage the object) is implemented using a badged
endpoint capability. The badge number is used to uniquely identify the instance
of the object. Since most expected implementations of the \refOS protocol will
use a similar method, the term \texttt{cap} (capability) may be used
interchangably with \texttt{object}. For example "window cap" means "a
capability which represents the access to a window object."

Note that in implementations, objects may be merged in some cases. For example, the process
server's \obj{liveness}, \obj{anon} or \obj{process} object caps may be merged into one for
simplification.

\begin{description}

	\item[\obj{process}] is an object, most likely maintained by the process
	server (i.e. \srv{procserv}), which represents a process. If a something
	has access to a process object, then it has the access to kill that
	process, call methods on behalf of that process...etc.
      
	\item[\obj{liveness}] is an object representing the identity of a process.
	If something has access to a liveness object of a process, then it can be
	notified of the process's death can ask for an ID to uniquely identify the
	process, but cannot kill the process or pretend to be it.
       
	\item[\obj{anon}] is an object representing the "address" of a server, used
	to establish a session connection to the server.
    
	\item[\obj{session}] is an object representing an active connection session
	to a server.  If something has access to a session object, then it can
	invoke methods on the server on behalf of the session client.
    
	\item[\obj{dataspace}] is an object that represents a dataspace. The
	dataspace itself may represent anything that may be modeled as a series of
	bytes including devices, RAM or files.  If something has access to a
	dataspace object, then it may read from / write to / execute it (depending
	on the dataspace permissions) by mapping a memory window to the dataspace,
	or close / delete the dataspace.
      
	\item[\obj{memwindow}] is an object representing an address space range
	(i.e. a memory window) segment in a process's virtual memory. If something has
	access to a memory window object, then it may map it to a dataspace, map a
	frame into it, and so forth.
    
\end{description}

% ----------------------------------------------------------------------
\section{Methods}

% ------------------
\subsection{Process Server Interface}

The process server interface aims to provide the abstraction for managing everything to do with a
process and its threads. This includes management of the process's virtual memory, memory window
creation / deletion, process creation / deletion, death notification of a client process, thread
management, and so on. Note that in implementations, the \cp{procserv}{session} may be the same
capability as \cp{procserv}{process} (in which case the process server would be connectionless).
It may also be shared with \cp{procserv}{anon} for simplification.

\begin{description}
	\item \pro{procserv}{session}{watch\_client(\cp{procserv}{liveness}, death\_notify\_ep)} 
	   {(Errorcode, death\_id, principle\_id)}   
    Authenticate a new client of a server against the \srv{procserv} and register for death
    notification.
    \begin{description}
        \item [procserv\_liveness\_C] The new client's liveness cap, which they have given to the
                server through session connection.
        \item [death\_notify\_ep] The async endpoint through which death notification will happen.
        \item [death\_id] The unique clientID that server will receive on death notification.
        \item [principle\_id] Optional, may be used for permission checking.
    \end{description}

    \item \pro{procserv}{session}{unwatch\_client(\cp{procserv}{liveness})}{(Errorcode)}
    Stop watching a client and remove death notifications about it.

	\item \pro{procserv}{session}{create\_mem\_window(base\_vaddr, size, permissions, flags)}
	   {(ErrorCode, \cp{procserv}{window})}
	Create a new memory window segment for the calling client.
	Note that clients may only create memory windows for their own address space, and alignment
	restrictions may apply here due to implementation and/or hardware restrictions. In the \refOS
    client environment, a valid memory window segment must be covering any virtual address ranges
    before any mapping can be done (including dataspace \& device frame mappings).
    \begin{description}
        \item [base\_vaddr] The window base address in the calling client's VSpace.
        \item [size] The size of the mem window.
        \item [permissions] Optional, read / write permission flags.
        \item [flags] Optional, extra flags bitmask (eg. cached / uncached window).
    \end{description}

    \item \pro{procserv}{session}{resize\_memory\_window(\cp{procserv}{window}, size)}
        {ErrorCode}
    Resize a memory window segment. Optional luxy feature, may be useful for implementing
    dynamic heap on clients.
    \begin{description}
        \item [size] The new size of the mem window to resize to.
    \end{description}

    \item \pro{procserv}{session}{delete\_memory\_window(\cp{procserv}{window})}
        {ErrorCode}
    Delete a memory window segment.

	\item \pro{procserv}{session}{register\_as\_pager(\cp{procserv}{window}, fault\_notify\_ep)} 
	   {(ErrorCode, window\_id)} 
    Register to receive faults for the presented window. The returned \ty{window\_id} is an integer
    used during notification in order for the pager to be able to identify which window faulted.
    \ty{window\_id} must be unique for each pager, although may also be implemented to be unique
    across the entire system.
    \begin{description}
        \item [fault\_notify\_ep] The async endpoint which the fault notifications will
                be sent through.
        \item [window\_id] The uniqie ID of the window, which will be used to identify which window
                faulted. The server most likely will have to book-keep this ID to handle faults
                correctly.
    \end{description}

    \item \pro{procserv}{session}{unregister\_as\_pager(\cp{procserv}{window})}{(ErrorCode)} 
    Unregister to stop being the pager for a client process's memory window.

	\item \pro{procserv}{session}{window\_map(\cp{procserv}{window}, window\_offset, src\_addr)}
	   {ErrorCode}
	This syscall is most commonly used in response to a prior fault notification from the
    process server. Maps the frame at the given VSpace into the client's faulted window, and
    then resolves the fault and resumes execution of the faulting client. Also may be used to
    eagerly map frames into clients before they VMfault there.
    \begin{description}
        \item [window\_offset] The offset into the window to map the frame into.
        \item [src\_addr] The address of the source frame in the calling client process's own
                VSpace; this address should contain a valid frame, and page-alignment restrictions
                may apply for this parameter.
    \end{description}
	
	\item \pro{procserv}{session}{new\_proc(name, params, block, priority)}{(ErrorCode, status)}
	Starts a new process, blocking or non-blocking.
    \begin{description}
        \item [name] The executable file name of the process to start.
        \item [params] The parameters to pass onto the new process.
        \item [block] Whether to block until the child process exits.
        \item [priority] The priority of the new child process.
        \item [status] Output, containint the exit status of the process. Only applicable if
                       blocking.
    \end{description}

    \item \pro{procserv}{session}{exit(status)}{(ErrorCode)}
    Exits and deletes the calling process.
    \begin{description}
        \item [status] The exit status of the calling client process.
    \end{description}

    \item \pro{procserv}{session}{clone(entry, stack, flags, args)}{(ErrorCode, thread\_id)}
    Starts a new thread, sharing the current client process' address space. The child
    thread will have the same priority as the parent process.
    \begin{description}
        \item [entry] The entry point vaddr of the new thread.
        \item [stack] The stack vaddr of the new thread.
        \item [flags] Any thread-related flags.
        \item [args] Thread arguments.
        \item [thread\_id] The thread ID of the cloned thread.
    \end{description}

\end{description}

% ------------------
\subsection{Server Connection Interface}

The server connection interface provides client session connection and disconnection. It may be a
good idea during implementation to extend this interface onto any other extra OS functioanlity that
is common across servers eg. debug ping, parameter buffer setup, notification buffer setup, and so
forth.

\begin{description}

    \item \pro{serv}{anon}{connect(\cp{procserv}{liveness})}{(ErrorCode, \cp{serv}{session})}
        Connect to a server and establish a session.

    \item \pro{serv}{session}{disconnect()}{(ErrorCode)}
        Disconnect from a server and delete session.

\end{description}

% ------------------
\subsection{Data Server Interface}

The data server interface provides the abstraction for management of dataspaces, including
creation, access, sharing, manipulation...etc.

\begin{description}
      
    \item \pro{dataserv}{session}{open(char *name, flags, mode, size)}
        {(ErrorCode, \cp{dataserv}{dataspace})}
    Opens a new dataspace at the dataspace server, which represents a series of bytes. Dataspace
    mapping methods such as datamap() and init\_data() directly or indirectly maps the contents of
    this dataspace into a memory window, after which the contents can be read from / written to.
    The concept of a dataspace in \refOS is similar to a file in UNIX; what a dataspace
    represents depends on the server that is implementing this interface.
    \begin{description}
        \item [name] The name of the dataspace to open.
        \item [flags] The read / write / create flags.
        \item [mode] The mode to create new file with, in the case that a new one is created.
        \item [size] The size of dataspace to open. Note that some data servers may ignore this.
    \end{description}

    \item \pro{dataserv}{session}{close(\cp{dataserv}{dataspace})}{ErrorCode}
    Close a dataspace belonging to the data server.

    \item \pro{dataserv}{session}{expand(\cp{dataserv}{dataspace}, size)}{ErrorCode}
    Resize and expand a given dataspace. Note that some dataspaces may not support this, as
    sometimes the size of a dataspaces makes no sense (eg. serial input).
    \begin{description}
        \item [size] The new size to expand dataspace to.
    \end{description}

    \item \pro{dataserv}{session}{datamap(\cp{dataserv}{dataspace}, \cp{procserv}{window}, offset)}
        {ErrorCode}
    Request that the data server back the specified window with the specified dataspace. Offset is the
    offset into the dataspace to be mapped to the start of the window. Note: the dataspace has to be
    provided by the session used to request the backing of the window.
    \begin{description}
        \item [procserv\_window\_C] Cap to the memory window to map the dataspace contents into.
        \item [offset] The offset in bytes from the beginning of the dataspace.
    \end{description}

    \item \pro{dataserv}{session}{dataunmap(\cp{procserv}{window})}
        {ErrorCode}
    Unmap the contents of the data from the given memory window.
    \begin{description}
        \item [procserv\_window\_C] Capability to window to unmap dataspace from.
    \end{description}

    \item \pro{dataserv}{session}{init\_data(\cp{dest}{dataspace}, \cp{dataserv}{dataspace},
        offset)}{ErrorCode}
    Initialises another dataspace by the contents of a source dataspace. The source dataspace is
    where the content is, and must originate from the invoked dataserver. Whether the destination
    dataspace and the source dataspace coming from the same dataserver is supported is up to the
    dataserver implementation; the dataserver documentation should be referred to. One example use
    case for this is a memory manager implementing the dataspace for RAM having a block of ram
    initialised by an external data source such as a file from a file server.
    \begin{description}
        \item [dest\_dataspace\_C] Another dataspace to be initialised with content from the 
                                 source dataspace.
        \item [dataserv\_dataspace\_C] The dataserver's own dataspace, where the content is.
        \item [offset] The content offset into the source dataspace.
    \end{description}

    \item \pro{dataserv}{session}{have\_data(\cp{dataserv}{dataspace}, fault\_ep)}
        {ErrorCode, data\_id}
    The local dataspace server calls another remote dataspace server to have the remote dataspace
    initalised by the contents of the local dataspace. The local dataspace must book-keep the given
    data\_id. The remote dataspace then will notify through the given endpoint asking for content
    initialisation, the dataID being given in the notification.
    \begin{description}
        \item [fault\_ep] The async endpoint to ask for content initialisation with.
        \item [data\_id] The remote endpoint's unique data ID number, used for notification.
    \end{description}

    \item \pro{dataserv}{session}{unhave\_data(\cp{dataserv}{dataspace})}{ErrorCode}
    Tell dataserver to stop providing content initialise data for its dataspace.

    \item \pro{dataserv}{session}{provide\_data(\cp{dataserv}{dataspace}, offset, content\_size)}
        {ErrorCode}  
    Gives the content from the local dataserver to the remote dataserver, in response to the
    remote server's earlier notification asking for content. The content is assumed to be in
    the set up parameter buffer. This call implicitly requires a parameter buffer to be
    set up, how this is done is up to the implementation.
    Note that even though the notification from the dataserver asking for content uses dataID to
    identify which dataspace, the reply, for security reasons, gives the actual dataspace cap.
    The dataID may be used securely iff the dataserver implementation supports per-client dataID
    checking, and a version of this method with dataID instead of a cap may be added.
    \begin{description}
        \item [offset] The offset into the remote dataspace to provide content for.
        \item [content\_size] The size of the content.
    \end{description}

    \item \pro{dataserv}{process}{datashare(\cp{dataserv}{dataspace})}{ErrorCode}
    Share a dataspace of a dataserver with another process. The exact implementation of this will be
    context-based. For example, a server sharing a parameter buffer may implement this function as
    \texttt{share\_param\_buffer(dataspace\_cap)}, implicitly stating the context, and implementing
    multiple share calls for more contexts. Although this method of stating the sharing context is
    strongly encouraged, the exact method of passing context is left up to the implementation. \\
    Note: There are a few assumptions made here.
    \begin{itemize}
    \item The dataspace is backed by somebody the receiving process trusts
        (eg. the process server).
    \item The dataspace is not revocable so the receiving process does
        not need to protect itself from an untrusted fault handler on that memory.
    \end{itemize}

    \item \pro{dataserv}{event}{pagefault\_notify(window\_id, offset, op)}{ErrorCode} \\
    Send a notification event to a data server indicating which window the fault occurred in,
    the offset within that window, and operation attempted (either Read or Write).

    \item \pro{dataserv}{event}{initdata\_notify(data\_id, offset)}{ErrorCode} \\
    Send a notification event to a data server indicating which dataspace needs its initial data.

    \item \pro{dataserv}{event}{death\_notify(death\_id)}{ErrorCode}
    Send a notification event to a data server indicating a client of
    theirs has died and the resources associated with them can be cleaned up.
  
\end{description}

% ------------------
\subsection{Name Server Interface}

The name server interface provides a common interface for servers to implement. We have chose a
simple hierachical distributed naming scheme here, where each server is responsible for a particular
prefix under another name server. This allows for simple distributed naming while allowing prefix
caching. A root name server must be defined from which naming starts; in our sample implementation
the process server acts as the root name server.

\begin{description}

    \item \pro{nameserv}{session}{register(name, \cp{dataserv}{anon})} {ErrorCode} \\
    Servers need to create an endpoint and call this on the root name server, in order for clients
    to be able to find this server and to connect to it. The anon cap will be given to clients
    looking for this server, and the client will make its connection call through the anon cap
    to establish a session.
    Re-registering will replace the current server anon cap.
    \begin{description}
        \item [name] The name to register under.
        \item [dataserv\_anon\_C] The anonymous endpoint to register with.
    \end{description}

    \item \pro{nameserv}{session}{unregister(name)} {ErrorCode} \\
    Unregister as a server under given name, so clients are no longer able to find us
    under that name. Existing anon caps will be invalidated.
    \begin{description}
        \item [name] The name to unregister for.
    \end{description}

    \item\pro{nameserv}{session}{resolve\_segment(path)}
        {(ErrorCode, \cp{dataserv}{anon}, resolved\_bytes)}
    Gives back an anon cap if the server is found. For details regarding the anon cap, refer to
    the documentation under register(). This function resolves a part of the name, returned
    as offset into the string that has been resolved. The rest of the path string may be
    resolved by another dataspace, and so forth, until we get to the end containing the endpoint
    to the server that actually contains the file. This allows for a simple hierachical
    namespace with distributed naming.
    \begin{description}
        \item [path] The path to resolve.
        \item [resolved\_bytes] The number of bytes resolved from start of path string.
    \end{description}

\end{description}

% ----------------------------------------------------------------------
\section{Server Components}

RefOS embraces the component-based multi-server OS design. The two main components in our design are
the process server and the file server. In our sample implementation we included an additoinal OS
server which is a dataserver responsible for basic device related tasks.

% ------------------
\subsection{Process Server}

The process server is the most trusted component in the system. It runs as the initial kernel thread
and does not depend on any other component (this avoids deadlock). The process server implementation
is single-threaded. The process server also implements the dataspace interface for anonymous memory,
and acts as the memory manager.

In our implementation, the process server implements the following interfaces:
\begin{itemize}
    \item Process server interface. (naming, mem. windows, processes..etc)
    \item Dataspace server interface. (For anonymous memory)
    \item Name server interface (In our implementation, acts as the root name server)
\end{itemize}

% ------------------
\subsection{File Server}

The file server is more trusted than clients, but less trusted than the process server (avoids
deadlock). In our current implementation the file server does not use a disk driver, and the actual
file contents are compiled into the file server executable itself using a pre-compiled archive.
The file server acts as the main data server in our system.

In our implementation, the file server implements the following interfaces.
\begin{itemize}
    \item Dataspace server interface. (for stored file data)
    \item Server connection interface. (for clients to connect to it)
\end{itemize}

% ------------------
\subsection{Console Server}

The Console server provides serial / EGA input \& output functionality, exposed through
the dataspace interface. It also provides terminal emulation for EGA screen output.

In our implementation, the Console server implements the following interfaces.

\begin{itemize}
    \item Dataspace server interface. (for serial input / output, and EGA screen devices)
    \item Server connection interface. (for clients to connect to it)
\end{itemize}

% ------------------
\subsection{Timer Server}

The Timer server provides timer gettime \& sleep functionality, exposed through
the dataspace interface.

In our implementation, the Timer server implements the following interfaces.

\begin{itemize}
    \item Dataspace server interface. (for timer devices)
    \item Server connection interface. (for clients to connect to it)
\end{itemize}


\chapter{Message Protocol}
\label{s:protocolip}
%
% Copyright 2014, NICTA
%
% This software may be distributed and modified according to the terms of
% the BSD 2-Clause license. Note that NO WARRANTY is provided.
% See "LICENSE_BSD2.txt" for details.
%
% @TAG(NICTA_BSD)
%

\section{Message Sequences}

This section will provide a series of example message sequences for common tasks, using the protocol
interfaces described.

\subsection{Server registration}

\autoref{f:servreg} illustrates the protocol a server uses to register itself under a name.

\begin{figure}[htb]
\begin{center}
\begin{msc}
msc {
    hscale = "1.5";

    R [ label ="Name Server"],
    S [ label = "Data Server"];
    
    S box S [ label = "create badged endpoint cap dataserv\_anon\_C"];
 
    S => R [ label = "nameserv\_session\_C.register(\"server\_name\", dataserv\_anon\_C)"];
    R box R [ label = "handle request, register under given name"];
    R >> S [ label = "ErrorCode"];
}
\end{msc}
\end{center}
\caption{Server registration.}
\label{f:servreg}
\end{figure}

\subsection{Establish a session to a server}


\autoref{f:estab} illustrates that to interact with a server, one must establish a session with it
through its anonymous endpoint cap. The end state after establishment should be an established
session to the server represented by a \obj{session} cap, with the server knowing the security
subject of the client, and optionally death notification set up so the server knows when the
client exits and can clean up book-keeping after it.

\begin{figure}[htb]
  \begin{center}
    \begin{msc}
      msc {
        hscale = "1.5";

        R [ label ="Process Server"],
        F [ label = "File Server"],
        C [ label = "Client"],
        N [ label ="Name Server"];

        R:>C [ label = "procserv\_liveness\_C"];
        C box C [label = "loop until path fully resolved:"];
        C=>N [ label = "nameserv\_session\_C.resolve\_segment(\"dataserver\")"];
        N>>C [ label = "(ErrorCode, dataserv\_anon\_C)"];
        C=>F [ label = "dataserv\_anon\_C.connect(procserv\_liveness\_C)"];
        F box F [label = "handle request, create client"];
        F=>R [ label = "procserv\_session\_C.watch\_client(procserv\_liveness\_C)"];
        R box R [ label = "set up death notify"];
        R>>F [ label = "(ErrorCode, death\_id, principle\_id)"];
        F box F [ label = "create session cap for new client"];
        F>>C [ label = "(dataserv\_session\_C)"];
      }
    \end{msc}
  \end{center}
  \caption{Example message sequence for finding \& establishing connection with a server.}
  \label{f:estab}
\end{figure}

Note the (optional) principle\_id is intended to be used for authentication assuming
ACL-like security model in the file system.

\pagebreak

\subsection{Opening and closing a Dataspace}

\autoref{f:openclose} show a sample where a client opens a dataspace which
could subsequently be used for mapping in a window. After use, it can be
closed.

\begin{figure}[htb]
  \begin{center}
    \begin{msc}
      msc {
        hscale = "1.5";

        D [ label = "Data Server"],
        C [ label = "Client"];

        C=>D [ label = "dataserv\_session\_C.open(name, flags, mode, size)"];
        D box D [label = "handle request, create dataspace"];
        D>>C [ label = "(ErrorCode, dataserv\_dataspace\_C)"];
        C box C [label="read/write/share/use dataspace"];
        C=>D [label = "dataserv\_session\_C.close(dataserv\_session\_C)"];
        D>>C [label = "ErrorCode" ];
      }
    \end{msc}
  \end{center}
  \caption{Sample protocol for opening and closing a dataspace.}
  \label{f:openclose}
\end{figure}


\subsection{Map Dataspace in a window}

\autoref{f:opendata} shows mapping an open dataspace into a given memory window.  The initialisation
part sets up the basic components for sample fault delegation in the system.  Note that the process
server is acting as the memory manager here in this example, as is in our simple implementation. The
paging could instead go through a separate memory manager process.

\begin{figure}[htb]
  \begin{center}
    \begin{msc}
      msc {
        hscale = "1.5";

        R [ label = "Process Server"],
        P [ label = "Data Server"],
        C [ label = "Client"];

        C abox P [ label = "Create a dataspace" ];
        C=>R [ label = "procserv\_session\_C.create\_mem\_window(base, size, perms, flags)"];
        R>>C [ label = "(ErrorCode, procserv\_window\_C)"];
        C=>P [ label = "dataserv\_session\_C.datamap(dataserv\_dataspace\_C, procserv\_window\_C, offset)"];
        P box P [label = "validate request"];
        P=>R [ label = "procserv\_session\_C.register\_as\_pager(procserv\_window\_C, dataserv\_notify\_C)"];
        R box R [ label = "book-keep pager information"];
        R>>P [ label = "(ErrorCode, window\_id)"];
        P>>C [ label = "ErrorCode"];
      }
    \end{msc}
  \end{center}
  \caption{Sample protocol for mapping a dataspace into a given window.}
  \label{f:opendata}
\end{figure}

\subsection{Page Fault}

\autoref{f:pagefault} show the page fault resolution protocol in the case of an external dataspace
server backing a window in the client's address space. In our implementation, the \seLf microkernel
sends the page fault notification to an endpoint cap pretending to be the faulting process.

\begin{figure}[h]
  \begin{center}
    \begin{msc}
      msc {
        hscale = "1.5";
        
        K [ label ="Kernel"],
        P [ label ="Process Server"],
        C [ label = "Client"],
        D [ label = "Data Server"];
       
        C box C [label = "VM fault into mapped dataspace"];
        K->P [label = "Page fault notification"];
        P box P [label = "transform fault into window\_id, offset"];
        P->D [ label = "dataserv\_event\_C.pagefault\_notify(window\_id,offset,op)"];
        D=>P [ label = "procserv\_session\_C.window\_map(window\_cap, window\_offset, src\_addr)"];
        P box P [ label = "handle request and find frame"];
        P:>C [ label = "map\_frame(frame)"];
        P>>D [ label = "ErrorCode"];
        P:>C [ label = "Resume execution"];
      }
    \end{msc}
  \end{center}
  \caption{Page fault protocol.}
  \label{f:pagefault}
\end{figure}

\clearpage
\subsection{Death notification}
\autoref{f:deathnotify} illustrates an example message sequence for death notification. Note that
it is assumed that the client is active and has already established connection with the data server
(see \autoref{f:estab}), and that the data server has since then set up death notification.

\begin{figure}[h]
\begin{center}
\begin{msc}
  msc {
    hscale = "1.5";

    C [ label = "Client"],
    P [ label = "Process Server"],
    D [ label = "Data Server"];
    
    C box C [ label = "Exits program" ];
    C => P [ label = "procserv\_session\_C.exit(status)" ];
    P box P [ label = "Delete client process" ];
    P >> D [ label = "dataserv\_event\_C.death\_notify(death\_id)" ];
    D box D [ label = "Delete client session book-keeping" ];
  }
\end{msc}
\end{center}
\caption{Example death notification.}
\label{f:deathnotify}
\end{figure}

\subsection{Dataspace Sharing}

\autoref{f:sharedata} is an example of how to set up a shared dataspace. Note that the assumption is
that the share is not revocable and that sharing is with someone who is trusted. So given two
clients, one must trust the other, or both must trust the other. Mutual distrust is not supported
yet simply as cap transfer to an untrusted entity requires more than a direct call.

Dataspace sharing is used to implement 

In this example, the shared dataspace is implemented by the process server (acting as the memory
manager), and is shared with a server. This is would be typical for a shared memory communication
buffer.

\begin{figure}[h]
  \begin{center}
    \begin{msc}
      msc {
        hscale = "1.5";

        R [ label ="Process Server"],
        C [ label = "Client"],
        D [ label = "Server"];
        
        C=>R [ label = "procserv\_session\_C.open(\"anon\", flags, mode, size)"];
        R>>C [ label = "(ErrorCode, procserv\_dataspace\_C)"];
        --- [ label = "Client maps the dataspace" ];
        C=>R [ label = "procserv\_process\_C.create\_mem\_window(base, size, perm, flags)"];
        R>>C [ label = "(ErrorCode, procserv\_window\_C)"];
        C=>R [ label = "procserv\_session\_C.datamap(procserv\_dataspace\_C, procserv\_window\_C, offset)"];
        R>>C [ label = "ErrorCode"];
        --- [ label = "Share to dataspace cap with server"];
        C=>D [ label = "serv\_session\_C.datashare(procserv\_dataspace\_C)"];
        D=>R [ label = "procserv\_process\_C.Create\_mem\_window(base, flags, mode, size)"];
        R>>D [ label = "(ErrorCode, procserv\_window2\_C)"];
        D=>R [ label = "procserv\_session\_C.datamap(procserv\_dataspace\_C, procserv\_window2\_C, offset)"];
        R>>D [ label = "ErrorCode"];
        D>>C [ label = "ErrorCode"];
        D abox C [ label = "Communicate using shared dataspace." ];
      }
    \end{msc}
  \end{center}
  \caption{Sample protocol for sharing an anonymous dataspace.}
  \label{f:sharedata}
\end{figure}

\clearpage

In the \autoref{f:sharedatathird} example, the dataspace is implemented not by the process server,
but instead an external data server. Thus, a third party is involved in this case. This is a more
general example of \autoref{f:sharedata}.

\begin{figure}[h]
  \begin{center}
    \begin{msc}
      msc {
        hscale = "1.5";

        R [ label ="Process Server"],
        D [ label = "Data Server"],
        CA [ label = "Client A"],
        CB [ label = "Client B"];
        
        CA=>D [ label = "dataserv\_session\_C.open(\"filename\", flags, mode, size)"];
        D>>CA [ label = "(ErrorCode, dataserv\_dataspace\_C)"];

        --- [ label = "Client A maps the dataspace" ];
        CA=>R [ label = "procserv\_session\_C.create\_mem\_window(base, size, perm, flags)"];
        R>>CA [ label = "(ErrorCode, procserv\_window\_C)"];
        CA=>D [ label = "dataserv\_session\_C.datamap(dataserv\_dataspace\_C, procserv\_window\_C, offset)"];
        D=>R [ label = "procserv\_session\_C.register\_as\_pager(procserv\_window\_C, dataserv\_notify\_C)"];
        R box R [ label = "Book-keep pager information"];
        R>>D [ label = "(ErrorCode, window\_id)"];
        D>>CA [ label = "ErrorCode"];

        --- [ label = "Now share the dataspace with client B"];
        CA=>CB [ label = "clientB\_session\_C.datashare(dataserv\_dataspace\_C)"];

        --- [ label = "Client B also maps the dataspace"];
        CB=>R [ label = "procserv\_session\_C.create\_mem\_window(base, size, perm, flags)"];
        R>>CB [ label = "(ErrorCode, procserv\_window2\_C)"];
        CB=>D [ label = "dataserv\_session\_C.datamap(dataserv\_dataspace\_C, procserv\_window2\_C, offset)"];
        D=>R [ label = "procserv\_session\_C.register\_as\_pager(procserv\_window2\_C, dataserv\_notify\_C)"];
        R box R [ label = "Book-keep pager information"];
        R>>D [ label = "(ErrorCode, window\_id)"];
        D>>CB [ label = "ErrorCode"];

        --- [ label = "Both clients have dataspace mapped"];
        CA abox CB [ label = "Communicate using shared dataspace." ];
      }
    \end{msc}
  \end{center}
  \caption{Sample protocol for sharing a dataspace provided by
    dataspace server between two clients.}
  \label{f:sharedatathird}
\end{figure}

\clearpage
\subsection{Content Initialise a Dataspace with Another Dataspace}

\autoref{f:initdata} shows creating a dataspace and setting it so that it is initialised by the
content of another. Out example initialises a RAM dataspace (with the process server acting as its
data server, as per our implementation) with the contents of an executable file from another
dataspace.

\begin{figure}[htb]
  \begin{center}
    \begin{msc}
      msc {
        hscale = "1.5";

        P [ label ="Process Server"],
        D [ label = "Data Server"],
        C [ label = "Client"];

        --- [ label = "Open the dataspace to be initialised"];
        C=>P [ label = "proc\_session\_C.open(\"anon\", flags, mode, size)"];
        P>>C [ label = "(ErrorCode, procserv\_dataspace\_C)"];

        --- [ label = "Open the data space to initialise with"];
        C=>D [ label = "dataserv\_session\_C.open(\"binary\", flags, mode, size)"];
        D>>C [ label = "(ErrorCode, dataserv\_dataspace\_C)"];

        --- [ label = "Initialise procserv dataspace with dataserv dataspace"];
        C=>D [ label = "dataserv\_session\_C.init\_data(procserv\_dataspace\_C, dataserv\_dataspace\_C, src\_offset)"];
        D box D [label = "Validate request and book-keep dataspace relationship"];
        D=>P [ label = "procserv\_process\_C.have\_data(procserv\_dataspace\_C, dataserv\_notify\_C)"];
        P box P [ label = "Book-keep pager information"];
        P>>D [ label = "(ErrorCode, data\_id)"];
        D>>C [ label = "ErrorCode"];

        --- [ label = "Now have dataspace initialised by external data server"];
        C=>P [ label = "procserv\_process\_C.create\_mem\_window(base, size, perm, flags)"];
        P>>C [ label = "(ErrorCode, procserv\_window\_C)"];
        C=>P [ label = "procserv\_session\_C.datamap(procserv\_dataspace\_C, procserv\_window\_C, offset)"];
        P>>C [ label = "ErrorCode"];

        C box C [ label = "Read from content-initialised mapped dataspace"];
      }
    \end{msc}
  \end{center}
  \caption{Sample protocol for creating a dataspace initialised with the contents of another.}
  \label{f:initdata}
\end{figure}

\clearpage
\subsection{Content Initialisated Page Fault}

\autoref{f:initfault} show the page fault resolution protocol in the case of an external dataspace
server providing the initial data for the dataspace. Note the similarity with \autoref{f:pagefault};
although different in concept, these two are similar in nature and may share a lot in
implementation. The subtle yet importance difference between the two is that for
\texttt{provide\_data} the content in the given frame is only read but not used; it is read and
copied to the receiving dataserver's own dataspace. This is in contrast to \texttt{window\_map}
which actually uses the given frame and maps the same physical frame into the client's address space
to resolve the fault.

\begin{figure}[htb]
  \begin{center}
    \begin{msc}
      msc {
        hscale = "1.5";
        
        K [ label = "Kernel"],
        P [ label ="Process Server"],
        C [ label = "Client"],
        D [ label = "Data Server"];
        
        C box C [label = "VM fault into mapped dataspace"];
        K->P [label = "Page fault notification"];
        P box P [label = "transform fault into window\_id, offset"];
        P->D [ label = "dataserv\_event\_C.initfault\_notify(data\_id, offset)"];
        D=>P [ label = "procserv\_session\_C.provide\_data(data\_id, src\_addr, offset)"];
        P box P [ label = "validate and find frame and copy data"];
        P:>C [ label = "map\_frame(frame)"];
        P>>D [ label = "ErrorCode"];
        P:>C [ label = "Resume execution"];
      }
    \end{msc}
  \end{center}
  \caption{Data content initialisation fault protocol.}
  \label{f:initfault}
\end{figure}

\subsection{ELF loading}

\autoref{f:elfloading} shows the elf loading protocol for bootstrapping a user process.

\begin{figure}[htb]
  \begin{center}
    \begin{msc}
      msc {
        hscale = "1.5";
        
        P [ label ="Process Server"],
        D [ label = "Data Server"],
        C [ label = "Client"];
        
        --- [ label = "User process starts at elf loading code in its own address space"];
        P abox C [ label = "Datamap ELF header file contents"];
        C box C [ label = "read ELF header"];

        --- [ label = "For each section (code/data):"];
        
        C => P [ label = "procserv\_session\_C.create\_mem\_window(base, size, perm, flags)"];
        P >> C [ label = "(ErrorCode, procserv\_window\_C)"];
        P abox C [ label = "Content initialise dataspace with code/data section content"];
        C => P [ label = "procserv\_session\_C.datamap(procserv\_dataspace\_C, procserv\_window\_C, offset)"];
        P >> C [ label = "ErrorCode"];

        P abox C [ label = "Create anonymous dataspace for zeroed section"];
        C => P [ label = "procserv\_session\_C.create\_mem\_window(base, size, perm, flags)"];
        P >> C [ label = "(ErrorCode, procserv\_window2\_C)"];
        P abox C [ label = "Datamap zeroed section into window"];

        --- [ label = "After all sections have been mapped:"];
        C box C [ label = "jump to program entry point"];
        C box C [ label = "set up userland environment"];
      }
    \end{msc}
  \end{center}
  \caption{Elf loading protocol.}
  \label{f:elfloading}
\end{figure}

\clearpage
\section{Data Structures}

This section will provide a series of example data structures for \refOS. These are only meant to be
a guideline, and will be presented in the C language. The data structures here aim to be the minimal
information such a structure should store, for a rough idea of what these objects represent.
Implementations may need to add extra members and references to these data structures to make
internal book-keeping easier.

\subsection {Memory window}

A window should book-keeps its state, which vspace it belongs to, where and how long the segment is,
and the state of its contents. A window my be unmapped, or mapped to a dataspace (and that dataspace
may be content initialised). In our implementation, the process server was also the anonymous
memory dataserver, so this window structure has extra book-keeping for internal anonymous dataspace
mapping state. Also, our implementation separated offset + size from a window.

\begin{verbatim}
struct w_window {
    vspace_t *vspace;
    vaddr_t offset;
    vaddr_t size;
    uint32_t permissions;
    bool cacheable;
    cspacepath_t capability;
    pager_t pager;
};
\end{verbatim}

\subsection {Process control block}

A process structure should book-keep its address space (capability space and virtual memory space),
its threads, and clients is it watching. In our implementation we also added parameter buffers,
notification ring buffers, process' OS system capabilities bitmask, parent, debug name, and so
forth.

\begin{verbatim}
struct proc_pcb {
    struct vs_vspace vspace;
    proc_tcb threads[];
    struct proc_watch_list clientWatchList;
};
\end{verbatim}

\subsection {Anonymous dataspace}

An anonymous dataspace would need to book keep its capability, list of frame pages it is
representing, and the content initialisation data. The content initialisation bitmask keeps which
pages in the dataspace have been initialised, and which still need to be initialised. Initialisation
happens lazily, on the first VM fault, so a waiting list of clients that are VM fault blocked
waiting for content to be initialised needs to be kept.

\begin{verbatim}
struct ram_dspace {
    cspacepath_t capability;
    vka_object_t *pages;
    uint32_t npages;
    cspacepath_t contentInitEP;
    uint32_t *contentInitBitmask;
    cvector_t contentInitWaitingList;
};
\end{verbatim}

\subsection {Archived file server dataspace}

This dataspace assumes that there is a char* storing the file data somewhere in the process server's
address space. In our implementation this came from a CPIO format archive linked into the process's
ELF sections.

\begin{verbatim}
struct fs_dataspace {
    seL4_CPtr dataspaceCap;
    seL4_Word permissions;
    char *fileData;
    size_t fileDataSize;
    content_init_info_t contentInitInfo;
};
\end{verbatim}

\subsection {Userland dataspace mapping}

This is the example data structure for a userland dataspace mapping. The window, dataspace, and
session are book-kept, along with the virtual address and size of the window segment.

\begin{verbatim}
typedef struct data_mapping {
    seL4_CPtr session;
    seL4_CPtr dataspace;
    seL4_CPtr window;
    char* vaddr;
    int windowSize;
    int dspaceSize;
} data_mapping_t;
\end{verbatim}


\chapter{Implementation Notes}
\label{s:implementation}
%
% Copyright 2014, NICTA
%
% This software may be distributed and modified according to the terms of
% the BSD 2-Clause license. Note that NO WARRANTY is provided.
% See "LICENSE_BSD2.txt" for details.
%
% @TAG(NICTA_BSD)
%

This section contains very brief overview of our implementation, which could be useful for anyone
implementing their own \refOS. For more details, refer to the Doxygen code domentation in our sample
implementation.

% ----------------------------------------------------------------------
\section{Bootstrapping}

\subsection*{Bootstrapping System Processes}
While implementing a multi-server OS, one quickly runs across the problem of trying to start
OS infrastructure with no OS infrastructure in place; The executables are started via a
bootstrapper app talking to the file server, but how do we start the file server itself?

In our implementation, we had the process server, being the root task and the most trusted
component, acting as a mini ELF loader in order to start the file server directly from its own CPIO
archive. There are several other ways to do this; one being the process server forks a copy of 
itself to act as the initial file server, or use a kernel which supposed multiple initial
tasks.

\subsection*{Bootstrapping Client Processes}
After the system processes have been started, the infrastructure is in place to start userland
processes. Starting userland processes require IPC communication, and in order to avoid the process
server directly communication with file servers (which causes backwards dependency), we instead
use an external bootstrapper. This bootstrapper may be a short-lived process or thread, which
simple reads the new userland process' ELF file and maps the sections.

In our implementation, we had a tiny bootstrapper called self-loader which was linked into a really
high reserved virtual address using a linker script, and it simply runs in the same address space as
the starting client process, setting up the ELF sections in its own address space before directly
jumping into the client ELF's entry point.

% ----------------------------------------------------------------------
\section{Communication}

We used the kernel communication mechanisms to implement server interface calls. We used a simple
Python script RPC stub generator which takes our RefOS interface (extended for extra functionality
and to make implementation easier) in a simple XML format and generates corresponding C stub code.
Parameters are marshalled and unmarshalled over the kernel thread IPC buffer. This isn't very
efficient, but requires the least amount of infrastructure. For long string parameters which 
are unsuited to the IPC buffer, and notifications, a shared dataspace buffer is used.

For async notifications, a shared dataspace with a one-way ring buffer protocol is used between the
sender and the reciever. A kernel async endpoint provides synchronisation for the buffer.


% ----------------------------------------------------------------------
\section{Anonymous Memory}

The process server recieves all the VM fault notifications from the kernel, and delegates it
accordingly. In our implementation, the process server acts as a data server to implement anonymous
memory. This simplifies the system, as page faults into content initialised memory only needs one
single delegation indirection, as opposed to two levels of delegation (first from process server to
memory server and then from memory server to data server).  alternatively, the process server
interface and memory server interface may be separated, and the task distributed to another process
if need be, and this memory server may act as the VM fault reciever.



\chapter{\label{s:concl}Conclusions}

We have presented a reference design for a multi-server OS designed for an \texttt{L4} based
microkernel. We have outlined a set of minimal design abstractions and methodologies which state
a standard protocol for common tasks, and provided sample message sequences.
On top of this, we provide a documented sample implementation of this protocol, in the form of a
small OS.

%\cleardoublepage
%\bibliographystyle{alpha}
%\bibliography{extra}

%%%%%%%%%%%% END REPORTS %%%%%%%%%%%%
\end{document}

% 
%%% Local Variables: %%%
%%% ispell-local-dictionary: "british" %%%
%%% End: %%%
